\documentclass[11pt]{amsart}
\usepackage{amssymb}
\usepackage[margin=1.15in]{geometry}
\usepackage[hidelinks]{hyperref}
\hypersetup{pdftitle={Free groups and free abelian groups as filtered colimits}}

% Stamp on every page: file name and compilation date.
\newcommand{\stamptime}{\the\numexpr(\time-30)/60\relax:\ifnum\numexpr\time-60*((\time-30)/60)\relax<10 0\fi\the\numexpr\time-60*((\time-30)/60)\relax}
\AddToHook{shipout/foreground}{%
 \put(0.5\paperwidth,-\paperheight+0.55cm){\makebox[0pt]{\tiny\ttfamily
  unit/doc/\jobname.tex \ --- \ compiled \today, \stamptime}}}

\newtheorem{theorem}{Theorem}[section]
\newtheorem{proposition}[theorem]{Proposition}
\newtheorem{lemma}[theorem]{Lemma}
\newtheorem{corollary}[theorem]{Corollary}
\theoremstyle{definition}
\newtheorem{definition}[theorem]{Definition}
\theoremstyle{remark}
\newtheorem{remark}[theorem]{Remark}
\newcommand{\code}[1]{\texttt{#1}}
\newcommand{\mere}[1]{\left\lVert #1\right\rVert}
\newcommand{\El}{\mathrm{El}}
\newcommand{\Z}{\mathbb{Z}}
\newcommand{\act}{\mathrm{act}}
\newcommand{\ev}{\mathrm{ev}}
\newcommand{\Link}{\mathrm{Link}}

\title[Free groups as filtered colimits]{Free groups and free abelian groups as filtered colimits}
\date{}

\begin{document}
\begin{abstract}
In univalent foundations with set quotients, the free group $F(X)$ and the
free abelian group $A(X)$ on an arbitrary set $X$ are constructed as
filtered colimits of the free structures on the finite stages
$f:[n]\to X$. On a finite set the free structures have decidable normal
forms (reduced words, integer vectors), and the colimit transfers what
these normal forms say to an arbitrary $X$. We obtain the universal
properties, the injectivity of $\eta:X\to F(X)$, and, for a locally small
set $E$, that $F(E)$ is locally small and that $\eta$ has small fibres.
The construction is written once, for any family of free structures on
finite stages satisfying a short list of properties; groups and abelian
groups are two instances. Everything is checked in Unit.
\end{abstract}
\maketitle

\section{Introduction}

We work in univalent type theory with function extensionality,
propositional truncation $\mere{-}$ and set quotients by proposition-valued
equivalence relations. Set quotients are effective: if $[a]=[b]$ then
$a\sim b$.

On a set with decidable equality the free group is the group of reduced
words, and equality of elements is decidable. On an arbitrary set $X$
reducedness still makes sense, but a word can no longer be reduced, since
we cannot decide whether a letter is followed by its inverse. The quotient
of all words by the equivalence relation generated by cancellation gives
the universal property, but the universal property alone does not give an
explicit criterion for equality in $F(X)$; the colimit gives one
directly.

We use instead that a set is a filtered colimit of finite sets. A
\emph{stage} is a finite set $[n]$ with a map $f:[n]\to X$, and $F(X)$ is
the colimit of the free groups $F[n]$. Two words over stages become equal
in $F(X)$ exactly when some further stage identifies them, and at a stage
equality is decidable. This answers questions such as: what is the fibre
of $\eta$ at $e_0-e_1+e_2$ in the free abelian group, when equality on
$E$ is not decidable? Equality in
$F(E)$ is a statement about a function between finite sets, equalities in
$E$ and an equality of normal forms, so it is small when $E$ is locally
small; and $\eta x=z$ forces $x$ to be one of the finitely many generators
occurring in a representative of $z$.

\smallskip\noindent\textbf{Results.} For any set $X$:
\begin{enumerate}
\item[(A)] $F(X)$ is a group with $\eta:X\to F(X)$, and every map $X\to G$
to a group extends uniquely to a homomorphism; likewise $A(X)$ for abelian
groups (Theorems~\ref{thm:group}, \ref{thm:up}).
\item[(B)] $\eta$ is injective (Theorem~\ref{thm:inj}).
\item[(C)] If $E$ is locally small, then $F(E)$ and $A(E)$ are locally
small and $\eta$ has small fibres (Theorem~\ref{thm:size}).
\end{enumerate}

\section{Finite stages}

The finite types we use are given by codes: $0$, $1$, and $a+b$ for codes
$a,b$, with $\El(0)=\emptyset$, $\El(1)=\{*\}$,
$\El(a+b)=\El(a)\sqcup\El(b)$. Sums of stages are then sums of codes, and
no arithmetic on $[n]$ is needed. Every $\El(c)$ has a decidable
equality, by induction on $c$. Two further facts are proved by
induction on the code:
\begin{itemize}
\item $\El(c)$ is inhabited or empty;
\item for a decidable predicate $p$ on $\El(c)$, either $p(i)$ for some
$i$, or $p(i)$ fails for every $i$.
\end{itemize}

\section{The free group on a discrete type}\label{sec:words}

Let $D$ be a type with decidable equality. A letter is $a$ or $a^{-1}$
for $a\in D$; a word is a list of letters; it is \emph{reduced} if no
letter is followed by its inverse. Reducedness is decidable.

The letter $x$ acts on words by
\[
\act_x(\mathrm{nil})=x,\qquad
\act_x(y\,w)=\begin{cases} w & \text{if } y=x^{-1},\\ x\,y\,w & \text{otherwise,}\end{cases}
\]
and a word $u=x_1\cdots x_n$ acts by $\ev_u=\act_{x_1}\circ\cdots\circ\act_{x_n}$.
Both preserve reduced words. Following van der Waerden, the group laws
come from the action.

\begin{lemma}\label{lem:words}
For reduced $s,t,w$, any letter $x$ and any words $u,r$:
\begin{enumerate}
\item[(C)] $\act_x(\act_{x^{-1}}w)=w$;
\item[(K)] $\ev_{\act_x r}(t)=\act_x(\ev_r t)$;
\item[(A)] $\ev_{\ev_u s}(t)=\ev_u(\ev_s t)$;
\item[(N)] $\ev_w(\mathrm{nil})=w$;
\item[(I)] $\ev_{u^{-1}}(\ev_u s)=s=\ev_u(\ev_{u^{-1}} s)$, where $u^{-1}$ is
the reversed word of inverses.
\end{enumerate}
\end{lemma}
\begin{proof}
(C) If $w=\mathrm{nil}$ or $w$ does not start with $x$, then
$\act_{x^{-1}}w=x^{-1}w$ and $\act_x$ cancels it. If $w=x\,w'$, then
$\act_{x^{-1}}w=w'$, and $w'$ does not start with $x^{-1}$ since $w$ is
reduced, so $\act_x w'=x\,w'$.
(K) If $r=y\,r'$ with $y=x^{-1}$, the left side is $\ev_{r'}t$ and the
right side is $\act_x\act_{x^{-1}}\ev_{r'}t$, equal by (C). Otherwise
$\act_x r=x\,r$ and both sides are $\act_x(\ev_r t)$.
(A) By induction on $u$, using (K).
(N) By induction on $w$: $\act_x w'=x\,w'$ when $x\,w'$ is reduced.
(I) By induction on $u$, using (C) and $\ev_{u v}=\ev_u\circ\ev_v$.
\end{proof}

\begin{proposition}
The reduced words form a group $F(D)$ with $u\cdot v=\ev_u(v)$, unit
$\mathrm{nil}$, inverse $\ev_{u^{-1}}(\mathrm{nil})$. Its generators
$a\mapsto a$ are distinct.
\end{proposition}
\begin{proof}
Associativity is (A), the left unit holds by definition, the right unit
is (N), and the inverse laws follow from (I) with (A) and (N). If the
one-letter words $a$ and $b$ are equal, $a$ occurs in $b$, so $a=b$.
\end{proof}

\begin{proposition}[universal property]\label{prop:up-words}
Let $G$ be a group and $f:D\to G$. Put $\Phi(x_1\cdots x_n)=f(x_1)\cdots
f(x_n)$, with $f(a^{-1})=f(a)^{-1}$. Then $\Phi$ is a homomorphism on
$F(D)$ extending $f$, and every homomorphism $\psi$ extending $f$ equals
$\Phi$.
\end{proposition}
\begin{proof}
$\Phi(\act_x r)=f(x)\Phi(r)$ because $f(x)f(x^{-1})=1$; by induction
$\Phi(\ev_u s)=\Phi(u)\Phi(s)$. For uniqueness, a reduced word $x\,w$ is
the product of the one-letter word $x$ and $w$, and
$\psi(a^{-1})=\psi(a)^{-1}$; induct on the word.
\end{proof}

\begin{proposition}[functoriality and support]\label{prop:push-words}
A map $s:D\to D'$ of discrete types gives the homomorphism
$s_*:F(D)\to F(D')$, $s_*(u)=$ the reduced form of $s$ applied to $u$.
Then $(t\circ s)_*=t_*\circ s_*$, $\mathrm{id}_*=\mathrm{id}$, and
$s_*$ depends only on the values of $s$. If a generator $b$ occurs in
$s_*u$, then $b=s(a)$ for a generator $a$ occurring in $u$. Moreover
$\Phi_f\circ s_*=\Phi_{f\circ s}$.
\end{proposition}
\begin{proof}
As (K): $\ev_{s(\act_x r)}(t)=\act_{s(x)}(\ev_{s r}t)$, since
$s(x^{-1})=s(x)^{-1}$. Hence $\ev_{s(\ev_u v)}(t)=\ev_{su}(\ev_{sv}t)$,
which gives both the homomorphism property (with (A)) and functoriality.
For the support: the letters of $\act_x w$ are among $x$ and the letters
of $w$, so the letters of a reduced form are among those of the word.
\end{proof}

\section{The free abelian group on a finite type}\label{sec:ab}

For a code $c$ let $\Z^c$ be the set of functions $\El(c)\to\Z$, an
abelian group under pointwise addition, with the generators $\delta_j$
($\delta_j(i)=1$ if $i=j$ and $0$ otherwise). We write $[P]\,x$ for $x$ if
the decidable statement $P$ holds and $0$ otherwise.

\subsection*{Finite sums} In an abelian group $G$, the sum of
$g:\El(c)\to G$ is defined by induction on the code:
$\sum_{0}g=0$, $\sum_{1}g=g(*)$,
$\sum_{a+b}g=\sum_a g\circ\mathrm{inl}+\sum_b g\circ\mathrm{inr}$.

\begin{lemma}\label{lem:sums}
\begin{enumerate}
\item $\sum_i (g(i)+h(i))=\sum_i g(i)+\sum_i h(i)$ and $\sum_i 0=0$.
\item (Fubini) $\sum_{i\in\El(c)}\sum_{j\in\El(d)}g(i,j)=\sum_{j}\sum_{i}g(i,j)$.
\item ($\delta$-rule) $\sum_i [i=w]\,h(i)=h(w)$ for $w\in\El(c)$.
\end{enumerate}
\end{lemma}
\begin{proof}
All three by induction on the code. For (1) on $a+b$ one uses
$(x+y)+(z+w)=(x+z)+(y+w)$. For (2) on $a+b$, the induction hypotheses and
(1). For (3) on $a+b$ with $w=\mathrm{inl}\,w'$: the terms on the $b$ side
are $0$, since $\mathrm{inr}\,j\ne\mathrm{inl}\,w'$, and the $a$ side is
$h(\mathrm{inl}\,w')$ by induction.
\end{proof}

\subsection*{Multiples} For $g\in G$ put $n\cdot g=g+\cdots+g$ ($n$
times) and $(-n-1)\cdot g=-((n+1)\cdot g)$.

\begin{lemma}\label{lem:mult}
\begin{enumerate}
\item $(x+y)\cdot g=x\cdot g+y\cdot g$, hence
$(\sum_i n_i)\cdot g=\sum_i n_i\cdot g$.
\item A homomorphism $\chi:\Z\to G$ is $n\mapsto n\cdot\chi(1)$.
\end{enumerate}
\end{lemma}
\begin{proof}
(1) First $(x+1)\cdot g=g+x\cdot g$ and $(x-1)\cdot g=-g+x\cdot g$, by
cases on $x$, using $-(a+b)=-a-b$. Addition on $\Z$ adds $y$ to $x$ by
$|y|$ successors or predecessors, and induction along this gives (1).
(2) $\chi(0)=0$, $\chi(x+1)=\chi(x)+\chi(1)$ and
$\chi(x-1)=\chi(x)-\chi(1)$; induct on $n$ and on $-n-1$.
\end{proof}

\subsection*{The groups $\Z^c$} For $m:\El(a)\to\El(b)$ put
\[
m_*u\,(z)=\sum_{m(i)=z}u(i)=\sum_{i}[m(i)=z]\,u(i).
\]

\begin{proposition}\label{prop:ab-push}
$m_*$ is a homomorphism; $(n\circ m)_*=n_*\circ m_*$,
$\mathrm{id}_*=\mathrm{id}$, $m_*$ depends only on the values of $m$, and
$m_*\delta_i=\delta_{m(i)}$. The generators are distinct. If
$\delta_j=m_*u$, then $j=m(i)$ for some $i$ (support). $\Z^0$ is trivial.
\end{proposition}
\begin{proof}
$m_*$ is additive by Lemma~\ref{lem:sums}(1). For functoriality,
\[
\begin{aligned}
n_*(m_*u)(w)&=\sum_z [n(z)=w]\sum_i [m(i)=z]\,u(i)
=\sum_i\sum_z [m(i)=z]\,[n(z)=w]\,u(i)\\
&=\sum_i [n(m(i))=w]\,u(i),
\end{aligned}
\]
by Fubini and the $\delta$-rule; $\mathrm{id}_*u=u$ and
$m_*\delta_i=\delta_{m(i)}$ are instances of the $\delta$-rule. If
$\delta_i=\delta_j$ then $\delta_i(j)=1$, so $i=j$. If $\delta_j=m_*u$,
then $1=\sum_i [m(i)=j]\,u(i)$; decide whether some $i$ has $m(i)=j$: if
none, every term is $0$ and so is the sum.
\end{proof}

\begin{proposition}[universal property]\label{prop:ab-up}
Let $G$ be an abelian group and $f:\El(c)\to G$. Then
$\Phi_f(u)=\sum_i u(i)\cdot f(i)$ is a homomorphism $\Z^c\to G$ with
$\Phi_f(\delta_i)=f(i)$; every homomorphism $\psi$ with
$\psi(\delta_i)=f(i)$ equals $\Phi_f$; and $\Phi_f\circ m_*=\Phi_{f\circ m}$.
\end{proposition}
\begin{proof}
$\Phi_f$ is additive by Lemmas~\ref{lem:sums}(1) and~\ref{lem:mult}(1);
$\Phi_f(\delta_i)=\sum_k [k=i]\,f(k)=f(i)$. Naturality:
\[
\Phi_f(m_*u)=\sum_z\sum_i \big([m(i)=z]\,u(i)\big)\cdot f(z)
=\sum_i\sum_z [m(i)=z]\,u(i)\cdot f(z)=\sum_i u(i)\cdot f(m(i)),
\]
by Lemma~\ref{lem:mult}(1), Fubini and the $\delta$-rule. Uniqueness, by
induction on the code. On $0$, $\Z^0$ is trivial and $\psi(0)=0$. On $1$,
$\Z^1=\Z$ and $\psi(n)=n\cdot\psi(1)=n\cdot f(*)$ by
Lemma~\ref{lem:mult}(2). On $a+b$, every $u$ is
$\mathrm{inl}_*u_a+\mathrm{inr}_*u_b$ with $u_a=u\circ\mathrm{inl}$,
$u_b=u\circ\mathrm{inr}$ (by the $\delta$-rule at each point), and
$\psi\circ\mathrm{inl}_*$ is a homomorphism $\Z^a\to G$ sending $\delta_i$
to $f(\mathrm{inl}\,i)$; by induction it is $\Phi_{f\circ\mathrm{inl}}$,
and likewise on $b$. So $\psi(u)=\Phi_{f\circ\mathrm{inl}}(u_a)+
\Phi_{f\circ\mathrm{inr}}(u_b)=\Phi_f(u)$.
\end{proof}

\section{The colimit}\label{sec:colimit}

The construction uses only the following data on the finite stages: for
each code $c$ a group $T(c)$ with generators $\El(c)\to T(c)$; for each
map $m$ a homomorphism $m_*$, functorial, depending only on the values of
$m$, sending generators to generators; distinct generators; the support
property; $T(c)$ trivial when $\El(c)$ is empty; and the universal
property of $T(c)$ for a class of target groups, natural in $c$
($\Phi_f\circ m_*=\Phi_{f\circ m}$). Sections~\ref{sec:words}
and~\ref{sec:ab} give two instances: $T(c)=F(\El c)$ for all groups and
$T(c)=\Z^c$ for abelian groups.

Fix a type $X$. A \emph{stage} is $s=(c,f)$ with $f:\El(c)\to X$; a map
of stages $(c,f)\to(d,g)$ is $m:\El(c)\to\El(d)$ with $g\circ m=f$.
Stages have sums $(c,f)+(d,g)=(c+d,[f,g])$.

\begin{lemma}[coequaliser]\label{lem:coeq}
Let $g:\El(n)\to X$ and let $(p_i,q_i)_{i\in\El(d)}$ be pairs in
$\El(n)$ with $g(p_i)=g(q_i)$. There is $r:\El(n)\to\El(n)$ with
$g\circ r=g$ and $r(p_i)=r(q_i)$ for all $i$.
\end{lemma}
\begin{proof}
Merge the pairs one at a time, starting from the identity. To merge
$(a,b)$ after $r_0$, put $r=\sigma\circ r_0$ with
$\sigma(z)=r_0(a)$ if $z=r_0(b)$, and $\sigma(z)=z$ otherwise. Then
$g\circ r=g$ because $g(r_0a)=g(a)=g(b)=g(r_0b)$; $r(a)=r_0(a)=r(b)$;
and pairs already merged by $r_0$ stay merged, as $r$ factors through
$r_0$.
\end{proof}

An \emph{element} is a stage $s$ with $u\in T(s)$. Two elements are
\emph{linked}, $\Link((s,u),(s',v))$, if there are a stage $t$ and maps
$m:s\to t$, $n:s'\to t$ with $m_*u=n_*v$.

\begin{lemma}
The relation $\mere{\Link}$ is an equivalence relation.
\end{lemma}
\begin{proof}
Reflexivity and symmetry are clear. Given links $(t_1,m_1,n_1)$ from $x$
to $y$ and $(t_2,m_2,n_2)$ from $y$ to $z$, go to $t_1+t_2$ and apply
Lemma~\ref{lem:coeq} to the pairs $(\mathrm{inl}\,n_1(j),\mathrm{inr}\,m_2(j))$,
$j$ in the stage of $y$; these have the same image in $X$. With the
resulting $r$, $(r\,\mathrm{inl}\,m_1)_*u=(r\,\mathrm{inl}\,n_1)_*v
=(r\,\mathrm{inr}\,m_2)_*v=(r\,\mathrm{inr}\,n_2)_*w$.
\end{proof}

Let $F(X)$ be the quotient of the elements by $\mere{\Link}$, and
$\eta(x)=[(1,x),\text{generator}]$.

\begin{theorem}\label{thm:group}
$F(X)$ is a group with
\[
[(s,u)]\cdot[(s',v)]=[(s+s',\ \mathrm{inl}_*u\cdot\mathrm{inr}_*v)],\quad
1=[(0,1)],\quad [(s,u)]^{-1}=[(s,u^{-1})].
\]
In the abelian instance it is commutative.
\end{theorem}
\begin{proof}
Links are compatible with the operations: from links $(t,m,n)$ and
$(t',m',n')$ one gets the link $(t+t',m+m',n+n')$ between the products,
since $(m+m')_*(\mathrm{inl}_*u\cdot\mathrm{inr}_*v)=\mathrm{inl}_*m_*u\cdot
\mathrm{inr}_*m'_*v$. So the operations descend to the quotient, and the
laws, which are equations in a set, may be checked on representatives.
There they hold up to a single map of stages: the associativity map
$(s+s')+s''\to s+(s'+s'')$, the maps $0+s\to s$ and $s+0\to s$, and the
folding map $s+s\to s$ together with $0\to s$ for the inverse laws. In
the abelian instance the swap $s+s'\to s'+s$ links $x\cdot y$ to
$y\cdot x$, since $T$ is commutative at each stage.
\end{proof}

\begin{theorem}[universal property]\label{thm:up}
Let $G$ be a target group and $f:X\to G$. There is a homomorphism
$\bar f:F(X)\to G$ with $\bar f\circ\eta=f$, and every homomorphism $\psi$
with $\psi\circ\eta=f$ equals $\bar f$.
\end{theorem}
\begin{proof}
Put $\bar f[(s,u)]=\Phi_{f\circ s}(u)$. Along a map of stages
$m:s\to t$, naturality gives $\Phi_{f\circ t}(m_*u)=\Phi_{f\circ t\circ m}(u)
=\Phi_{f\circ s}(u)$, so $\bar f$ is constant on links and descends. It is
a homomorphism because each $\Phi$ is, and $\bar f(\eta x)=f(x)$. For
uniqueness, restrict $\psi$ to a stage $s$: $u\mapsto\psi[(s,u)]$ is a
homomorphism on $T(s)$ (the folding map $s+s\to s$ links
$(s,u)\cdot(s,v)$ to $(s,uv)$), sending the generator $i$ to
$\psi(\eta(s\,i))=f(s\,i)$ (the map $1\to s$ picking $i$). By the
universal property of $T(s)$ it is $\Phi_{f\circ s}$.
\end{proof}

\section{The unit is injective}

\begin{theorem}\label{thm:inj}
If $X$ is a set, $\eta:X\to F(X)$ is injective.
\end{theorem}
\begin{proof}
Let $\eta a=\eta b$. By effectiveness there merely is a link: a stage
$t$ and $m,n$ from the one-point stages $a$ and $b$ with
$m_*(\text{gen})=n_*(\text{gen})$, i.e.\ the generators $m(*)$ and $n(*)$
of $T(t)$ are equal. Generators are distinct, so $m(*)=n(*)$, and
$a=t(m(*))=t(n(*))=b$. Since $X$ is a set, $a=b$ is a proposition and the
link may be used although it exists only merely.
\end{proof}

\section{Size}

Let $E$ be a set in a universe $U'$ which is \emph{locally small} with
respect to a smaller universe $U$: each identity type $a=_E b$ is
equivalent to a type in $U$. A type is \emph{small} if it is equivalent
to a type in $U$. The finite stages, and the groups $T(c)$, are in $U$.

\begin{lemma}\label{lem:smalllink}
For elements $x=(s,u)$ and $y=(s',v)$, with $s+s'=(c,g)$,
\[
[x]=[y]\iff \mere{\textstyle\sum_{r:\El(c)\to\El(c)}\ \big(\prod_z g(r z)=g(z)\big)\times
(r\,\mathrm{inl})_*u=(r\,\mathrm{inr})_*v}.
\]
The right side is small.
\end{lemma}
\begin{proof}
From right to left, $(s+s',r\,\mathrm{inl},r\,\mathrm{inr})$ is a link.
From left to right, take a link $(t,m,n)$ and let $\rho=[m,n]:\El(c)\to
t$. If $\El(c)$ is empty, $T(c)$ is trivial and $r=\mathrm{id}$ works.
Otherwise choose a point $d$ and, for $z\in t$, let $\sigma(z)$ be some
$i$ with $\rho(i)=z$ if there is one (decided on the finite type
$\El(c)$), and $d$ otherwise. Then $r=\sigma\circ\rho$ satisfies
$\rho\circ r=\rho$, hence $g\circ r=g$, and
$(r\,\mathrm{inl})_*u=\sigma_*m_*u=\sigma_*n_*v=(r\,\mathrm{inr})_*v$.
The right side involves a function between finite types, equalities in
$E$, which are small, and an equality in $T(c)$, which is in $U$.
\end{proof}

\begin{lemma}\label{lem:smallprop}
Being small is a proposition.
\end{lemma}
\begin{proof}
Given $(Y,e)$ and $(Y',e')$ with $e:Y\simeq B$, $e':Y'\simeq B$, the
equivalence $d=e'^{-1}\circ e:Y\simeq Y'$ satisfies $e'\circ d=e$. By
univalence for $U$ (induction on the equivalence $d$) it suffices to
treat $Y'=Y$ and $d=\mathrm{id}$, where $e'=e$ because equivalences with
the same underlying map are equal (being an equivalence is a
proposition).
\end{proof}

\begin{theorem}\label{thm:size}
If $E$ is a locally small set, then $F(E)$ is locally small and the
fibres of $\eta$ are small: for $z=[(s,u)]$,
\[
\big(\textstyle\sum_{a:E}\eta a=z\big)\ \simeq\ \mere{\textstyle\sum_{i:\El(s)}\eta(s\,i)=z}.
\]
The same holds for $A(E)$.
\end{theorem}
\begin{proof}
On representatives, local smallness is Lemma~\ref{lem:smalllink}. For
the fibre: if $\eta a=z$, a link from the one-point stage $a$ to $(s,u)$
gives a stage $t$ and maps $m,n$ with $\text{gen}_{m(*)}=n_*u$. By the
support property $m(*)=n(i)$ for some $i\in\El(s)$, and
$a=t(m(*))=t(n(i))=s(i)$. The fibre is a proposition by
Theorem~\ref{thm:inj}, and the right side is small by
Lemma~\ref{lem:smalllink}. Every element of $F(E)$ is merely a
representative, and by Lemma~\ref{lem:smallprop} the statements are
propositions, so they hold at every element.
\end{proof}


\section{The free abelian group}\label{sec:abfree}

Take $T(c)=\Z^c$ in Section~\ref{sec:colimit} and let $A(X)$ be the
resulting colimit. An element of $A(X)$ is represented by a stage
$f:\El(c)\to X$ and an integer vector $u\in\Z^c$, read as the formal sum
$\sum_i u(i)\,f(i)$; the unit is $\eta(x)=[(1,x),1]$.

\begin{theorem}
$A(X)$ is an abelian group, and it is the free abelian group on $X$: every
map $X\to G$ to an abelian group extends uniquely to a homomorphism
$A(X)\to G$. If $X$ is a set, $\eta$ is injective; if $X$ is a locally
small set, $A(X)$ is locally small and $\eta$ has small fibres.
\end{theorem}
\begin{proof}
Commutativity is in Theorem~\ref{thm:group} (the swap of stages); then
$A(X)$ is a target, and Theorems~\ref{thm:up}, \ref{thm:inj}
and~\ref{thm:size} apply with Propositions~\ref{prop:ab-push}
and~\ref{prop:ab-up}.
\end{proof}

Unfolding Lemma~\ref{lem:smalllink} for this instance describes equality
in $A(E)$ by integer vectors. An endomap $r$ of $\El(c+c')$ with
$[f,f']\circ r=[f,f']$ identifies formal variables which have the same
value in $E$; the condition $(r\,\mathrm{inl})_*u=(r\,\mathrm{inr})_*v$ says
that for every $w$,
\[
\sum_{r(\mathrm{inl}\,i)=w}u(i)=\sum_{r(\mathrm{inr}\,j)=w}v(j).
\]

\begin{corollary}\label{cor:abeq}
$\sum_i u(i)\,f(i)=\sum_j v(j)\,f'(j)$ in $A(E)$ if and only if there
merely is a partition of the formal variables of the two sides into
classes on which $[f,f']$ is constant, such that in each class the
coefficients of the left side and of the right side have the same total.
\end{corollary}
\begin{proof}
A map $r$ as above gives the partition into its fibres, with the class
of $w$ having the totals of the fibre over $r(w)$. Conversely, a partition
(a decidable equivalence relation) gives $r$, sending each variable to the
first element of its class in the order of the code; this choice depends
only on the class, so the fibres of $r$ are the classes.
\end{proof}

\begin{corollary}\label{cor:abfinite}
Equality $\sum_i u(i)\,f(i)=\sum_j v(j)\,f'(j)$ in $A(E)$ is a finite
disjunction of conjunctions of equalities in $E$: it holds iff merely, for
one of the finitely many maps $r:\El(c+c')\to\El(c+c')$ whose fibre sums
agree (a decidable condition), $[f,f'](r(z))=[f,f'](z)$ for every $z$.
\end{corollary}
\begin{proof}
The maps $\El(a)\to\El(b)$ are coded by a finite type $\El(\mathrm{fn}(a,b))$,
with $\mathrm{fn}(0,b)=1$, $\mathrm{fn}(1,b)=b$ and
$\mathrm{fn}(a+a',b)=\mathrm{fn}(a,b)\times\mathrm{fn}(a',b)$, by a
decoding which is surjective up to pointwise equality; the condition on
the fibre sums is a finite conjunction of equalities of integers.
\end{proof}

\smallskip\noindent\textbf{Examples.} Let $e_0,e_1,e_2\in E$.
\begin{enumerate}
\item $e_0-e_1+2e_2=e_0+e_1$ holds iff $e_1=e_2$. Write the left side as
$X_0-X_1+2X_2$ and the right side as $Y_0+Y_1$, with $X_0,Y_0\mapsto e_0$
and $X_1,Y_1\mapsto e_1$, $X_2\mapsto e_2$. Counting the right side
negatively, the coefficients are $1,-1,2,-1,-1$, and every class must
have total $0$. The only negative coefficients are the three $-1$'s, so the
class of $X_2$ contains two of $X_1,Y_0,Y_1$ (or all other variables);
in either case it contains $X_1$ or $Y_1$, whose value is $e_1$, so
$e_1=e_2$. Conversely, if $e_1=e_2$ the classes $\{X_0,Y_0\}$ and
$\{X_1,X_2,Y_1\}$ work.
\item The fibre of $\eta$ at $e_0-e_1+e_2$ is $\mere{e_0=e_1\vee e_1=e_2}$.
By Theorem~\ref{thm:size} an element in the fibre is one of
$e_0,e_1,e_2$, and $\eta(e_i)=e_0-e_1+e_2$ asks for a partition of
$Y,X_0,X_1,X_2$ with coefficients $-1,1,-1,1$ and all totals $0$, with
$Y\mapsto e_i$. For $i=0$ the possibilities need $e_1=e_2$ (classes
$\{Y,X_0\},\{X_1,X_2\}$) or $e_0=e_1=e_2$; for $i=2$, $e_0=e_1$ or
$e_0=e_1=e_2$; for $i=1$, $e_0=e_1=e_2$.
\item The fibre of $\eta$ at $e_0-e_1+2e_2$ is empty. Here no partition is
needed: the augmentation $A(E)\to\Z$, the homomorphism sending every
generator to $1$, sends $e_0-e_1+2e_2$ to $2$ and every $\eta(a)$ to $1$.
\end{enumerate}

\section{What is formalised}

Everything above is checked in Unit, in \code{unit/sag/freegroup}. The
development has the parameters of \code{sag/}: univalence and replacement
(images and their surjectivity). Function extensionality, propositional
truncation and effective set quotients are derived from them. There are no
new axioms and no higher inductive types.

\smallskip
\begin{center}
\begin{tabular}{ll}
Section & Files \\ \hline
2 & \code{Decisions}, \code{Codes} \\
3 & \code{Words}, \code{Functor}, \code{Universal} \\
4 & \code{FinSum}, \code{Multiples}, \code{AbelianInstance} \\
5 & \code{Interface}, \code{Stages}, \code{Colimit}, \code{FreeGroup}, \code{Lift}, \code{FreeInstance} \\
6 & \code{Injective} \\
7 & \code{Identification}, \code{SmallProp}, \code{Small} \\
8 & \code{AbelianInstance}, \code{Abelian}, \code{AbelianEquality}, \code{AbelianExamples} \\
\end{tabular}
\end{center}
\smallskip

\noindent The list of properties at the beginning of
Section~\ref{sec:colimit} is the record \code{FinFree}; the colimit files
are written for an arbitrary \code{FinFree}, and \code{FreeInstance.free}
and \code{AbelianInstance.abelian} are the two instances. The final
statements are in \code{Theorems.seg}, with explicit types:
\code{free-group}, \code{free-lift-mul}, \code{free-lift-eta},
\code{free-lift-unique}, \code{free-eta-inj}, \code{free-locally-small},
\code{free-small-fibres}, and the same with prefix \code{ab-} for abelian
groups, together with \code{ab-comm}. The universal property is stated as
existence of the homomorphism and pointwise uniqueness.
For $A(X)$, \code{ab-eq-identification}, \code{ab-eq-partition} and
\code{ab-eq-finite} are the description by maps $r$ with equal fibre sums,
Corollary~\ref{cor:abeq} and Corollary~\ref{cor:abfinite};
\code{ab-example1}, \code{ab-example2}, \code{ab-example3} are the
three examples, for a set $E$. In the formal proofs of the first two
examples the partitions are not enumerated: from a map $r$, the fibre sums
at one or two points are integer expressions in the decidable memberships
$[r(Z)=w]$, and a case analysis on these memberships shows which variables
share a class. The third is the augmentation, an instance of the universal
property with $G=\Z$.

Two universe levels are used: the finite stages live in $U_l$ and the
generators in $U_k$ with $l\le k$, so that for $E$ in $U_{l+1}$ locally
small the groups $T(c)$ stay in $U_l$. The integers, with the arithmetic of $\mathbb N$ they
use, are those of \code{sag/stoneseq}; that being an equivalence
is a proposition is \code{sag/em/EquivProp}.

The development consists of 24 new files (2777 lines) and links 13
existing files. Checking it takes 0.22\,s (1255 declarations); no
declaration unfolds more than 42 definitions (the most are the explicit
fibre sums of the examples).

\end{document}
